\documentclass[14pt]{extarticle}
\usepackage{fontspec}
\usepackage{polyglossia} % use instead of babel
\usepackage{pgfplots}
\usepackage{float}
\usepackage{graphicx}
\usepackage{amsmath}
\usepackage{amsfonts}
\setdefaultlanguage{hebrew}
\setotherlanguage{english}
\newfontfamily\hebrewfont[Script=Hebrew]{Arial}

\usepackage[margin=1in]{geometry}

% Reset equation counter at the start of each section
\usepackage{chngcntr}
\usepackage{etoolbox}
\counterwithin*{equation}{section}
\preto{\section}{\setcounter{equation}{0}}

% Left-aligned equation environment
\newenvironment{lefteqs}{%
    \begin{english}%
    \begin{flushleft}%
    \renewcommand{\arraystretch}{1.3}%
    $\begin{array}{@{}l@{}}%
}{%
    \end{array}$%
    \end{flushleft}%
    \end{english}%
}

\begin{document}

\begin{center}
    {\LARGE \textbf{מעבדה 6}}\\
    {\textbf{מדידת מהירות האור}}
\end{center}

\begin{itemize}
    \item מגישים
    \begin{itemize}
        \item שי רואימי - 318986270
    \end{itemize}
    \item תאריך ביצוע הניסוי:
\end{itemize}

\newpage
\section*{ניסוי 1 - מדידת מהירות האור באוויר (שיטה 1)}

\subsection*{מטרת הניסוי}
מדידת מהירות אור בעזרת מדידת הפרש המרחק שעוברת קרן לייזר, כאשר הפרש הפאזה קבוע.

\subsection*{רקע תיאורטי}
אור הוא גל אלקטרומגנטי.
גל המתקדם לאורך ציר $x$ מתואר על ידי פונקציית הגל:
\begin{equation}
    \psi(x,t) = A \sin\left[ 2\pi \left( f t - \frac{x}{\lambda} \right) \right]
\end{equation}

כאשר $A$ המשרעת, $f$ התדירות, $\lambda$ אורך הגל, $x$ המרחק מהמקור, ו$t$ הזמן.

הביטוי בסוגריים הוא המופע (הפאזה) של הגל:
\begin{equation}
    \varphi(x,t) = 2\pi \left( f t - \frac{x}{\lambda} \right)
\end{equation}

הפרש המופע בין שתי נקודות הנמצאות על ציר התקדמות הגל במרחק $\Delta x$ הוא:
\begin{equation}
    \Delta \varphi = \varphi_1 - \varphi_2 = 2\pi \frac{x_2 - x_1}{\lambda} = 2\pi \frac{\Delta x}{\lambda}
\end{equation}

הקשר בין אורך הגל, התדירות ומהירות ההתפשטות בתווך הוא:
\begin{equation}
    \lambda = \frac{v}{f}
\end{equation}

הצבת (4) ב-(3) ובידוד $v$ נותנת את נוסחת המדידה:
\begin{equation}
    v = 2\pi \cdot \frac{\Delta x}{\Delta \varphi} \cdot f
\end{equation}

מכיוון שמד מהירות האור מציג את הפרש המופע במעלות, המרה לרדיאנים תתן לנו:
\begin{equation}
    v = 360 \cdot \frac{\Delta x}{\Delta \varphi} \cdot f
\end{equation}
כאשר $\Delta \varphi$ נמדד במעלות.

במערכת שלנו, נשים לב כי אנחנו לא מודדים את הגל האופטי עצמו (בעל תדירות מסדר גודל של $10^14 Hz$)
אלא את האות המאופנן בעל תדר $50 MHz$.
האות המאפנן מתקדם במרחב באותה מהירות של גל האור הנושא אותו ולכן מדידת זמן המעבר שלו היא מדידה של מהירות האור.

בנוסף, נציין שעבור מרחק הזזה של המראה $d$ במהלך הניסוי, מכיוון שקרן הלייזר יוצאת מהמכשיר ומוחזרת מהמראה חזרה לגלאי
אז הפרש הדרכים הוא כפול:

\begin{equation}
    \Delta x = 2d
\end{equation}

בניסוי זה נקבע $\Delta \varphi = 180^\circ$ ונמדד $d$ המתאים לו. הצבת (7) ב-(6) עם $\Delta \varphi = 180^\circ$ נותנת:
\begin{equation}
    c = 360 \cdot \frac{2d}{180} \cdot f = 4 d f
\end{equation}

\subsection*{מערכת המדידה ומהלך הניסוי}

הרכבנו על מסילה בעלת סרגל מטרי את מד מהירות האור ואת המראה המחזירה.
מד מהירות האור מכיל את מקור האור (דיודת \begin{english}LED\end{english} אדומה שמשרעתה מאופננת בתדר $50MHz$),
את הגלאי (פוטודיודה) ואת לוח התצוגה הדיגיטלי.

המכשיר פעל יותר מ-10 דקות לפני הגעתנו לעמדת הניסוי, ולכן התחלנו מיד בהכנת הניסוי.
יישרנו את המכשיר ואת המראה על המסילה כך שהקרן היוצאת מהמכשיר חזרה לפתח (ניתן לראות את נקודת קרן הלייזר, ועל פי זה לכוון).


לחצנו על כפתור \begin{english}"mode"\end{english} עד שהמכשיר עבר למצב תצוגת הפרש מופע במעלות.

מיקמנו את המראה קרוב ככל הניתן למקור האור, אפסנו את תצוגת המופע באמצעות כפתור \begin{english}"Calibration"\end{english}
ורשמנו את מיקום המראה.

לאחר מכן הרחקנו את המראה לאורך המסילה עד שקיבלנו הפרש פאזה של $180^\circ$, ורשמנו שוב את מיקום המראה.

חזרנו על הפעולה הזאת ארבע פעמים נוספות, עם מיקום התחלתי שונה.

על מנת להעריך את שגיאת המדידה, ביצענו גם מדידה של כמה ניתן להזיז את המראה מבלי לקבל שינוי בקריאת המופע במד מהירות האור.

\subsection*{תוצאות וניתוח}

תוצאות המדידה, כאשר $d$ הוא הפרש קריאות המיקום:

\begin{table}[H]
    \centering
    \renewcommand{\arraystretch}{1.5}
    \begin{tabular}{|c|c|c|}
        \hline
        \text{[cm] מיקום התחלתי} & \text{[cm] מיקום ב-$180^\circ$} & \text{[cm] $d$}\\ \hline
        0  & 153.5 & 153.5 \\ \hline
        3  & 155.5 & 152.5 \\ \hline
        5  & 156.6 & 151.6 \\ \hline
        7  & 158.1 & 151.1 \\ \hline
        10 & 163.6 & 153.6 \\ \hline
    \end{tabular}
    \caption{מרחק הזזת המראה הדרוש להפרש מופע של $180^\circ$}
    \label{tab:l6_e1}
\end{table}

הערך הטוב ביותר עבור $d$ הוא ממוצע חמש המדידות (נוסחה (1) בפרק השגיאות):

\begin{lefteqs}
    \bar{d} = \frac{1}{N}\sum_{i=1}^{N} d_i = \frac{153.5 + 151.6 + 153.6 + 152.5 + 151.1}{5} = 152.46\,cm
\end{lefteqs}

חמש המדידות בוצעו במיקומי מוצא שונים, ולכן הפיזור ביניהן מייצג את
אי-הוודאות האקראית של תהליך המדידה. לפי נוסחה (2) בפרק השגיאות השגיאה הסטטיסטית היא:
\begin{lefteqs}
    \Delta_s d = \sqrt{\frac{\sum_{i=1}^{N}\left(d_i - \bar{d}\right)^2}{N}} \\
    \Delta_s d = \sqrt{\frac{1.04^2 + (-0.86)^2 + 1.14^2 + 0.04^2 + (-1.36)^2}{5}} = 1.00\,cm
\end{lefteqs}

כאשר $\Delta_s d$ מייצג את השגיאה הסטטיסטית.

לגבי שגיאת המדידה, ישנן 2 שגיאות מדידות.
האחת היא שגיאת המדידה בסרגל, והשנייה רזולוציית תצוגת הפרש המופע במכשיר מד מהירות האור.
ראינו כי ניתן להזיז את המראה $7mm$ מבלי שהקריאה במד מהירות האור, כלומר לכל קריאת מופע מהצג
השגיאה היא מחצית התחום: $3.5mm$

הגודל $d$ הוא הפרש של שתי קריאות ולכן:

\begin{lefteqs}
    \Delta_m d = 0.35 + 0.35 = 0.7\,cm
\end{lefteqs}

כאשר $\Delta_m d$ מייצג את שגיאת המדידה.

שגיאת המדידה בסרגל היא $0.5mm$ וקטנה בסדר גודל מ $\Delta_m d$ ולכן נבלעת בה.

יש לציין כי למכשיר הדיוק בתצוגת המופע הוא $1^{\circ}$, כאשר ע"פ משוואת (6) ו-(7) ברקע התיאורטי
$1^\circ$ שקול להזזת מראה של:
\begin{lefteqs}
    \frac{c}{2 \cdot 360 \cdot f} = \frac{2.998 \cdot 10^8}{720 \cdot 50 \cdot 10^6} = 8.3\,mm
\end{lefteqs}
בהתאמה טובה ל-$7mm$ שנמדדו בפועל.

את השגיאה הכוללת נחשב לפי הנוסחא:
\begin{lefteqs}
    \Delta d = \sqrt{\left(\Delta_m d\right)^2 + \left(\Delta_s d\right)^2} = \sqrt{0.7^2 + 1.00^2} = 1.22\,cm
\end{lefteqs}

לפי (7), הפרש הדרכים ושגיאתו הם:
\begin{lefteqs}
    \Delta x = 2\bar{d} = 2 \cdot 152.46 = 304.92\,cm \\
    \delta(\Delta x) = 2 \Delta d = 2.44\,cm
\end{lefteqs}

כעת נחשב את מהירות האור לפי (8):
\begin{lefteqs}
    c = 4 \bar{d} f = 4 \cdot 1.5246 \cdot 50 \cdot 10^6 = 3.049 \cdot 10^8\,\frac{m}{s}
\end{lefteqs}

$c$ הוא מכפלה ומנה של הגדלים הנמדדים $\Delta x$, $\Delta \varphi$ ו-$f$.
לגודל מסוג זה השגיאה היחסית הכוללת היא סכום השגיאות היחסיות של המרכיבים (נוסחה (8) בתדריך השגיאות, ומצויינת גם במדריך הניסוי):
\begin{lefteqs}
    \frac{\delta c}{c} = \frac{\delta(\Delta \varphi)}{\Delta \varphi} + \frac{\delta(\Delta x)}{\Delta x} + \frac{\delta f}{f}
\end{lefteqs}
כאשר $\delta(\Delta \varphi) = 1^\circ$ ו- $\frac{\delta f}{f} = 5\%$ לפי נתוני היצרן.
נציב:
\begin{lefteqs}
    \frac{\delta c}{c} = \frac{1}{180} + \frac{2.44}{304.92} + 0.05 = 0.0056 + 0.0080 + 0.05 = 0.0635
\end{lefteqs}
כלומר שגיאה יחסית של $6.4\%$, ולכן:
\begin{lefteqs}
    \delta c = 0.0635 \cdot 3.049 \cdot 10^8 = 0.19 \cdot 10^8\,\frac{m}{s}
\end{lefteqs}

התוצאה הסופית:
\begin{lefteqs}
    c = (3.05 \pm 0.19) \cdot 10^8\,\frac{m}{s}
\end{lefteqs}

\subsection*{דיון ומסקנות}

התוצאה שהתקבלה למהירות האור באוויר היא:
\begin{lefteqs}
    c = (3.05 \pm 0.19) \cdot 10^8\,\frac{m}{s}
\end{lefteqs}

הערך הידוע של מהירות האור בריק הוא $2.998 \cdot 10^8\,\frac{m}{s}$.

ע"פ הספרות מקדם השבירה באוויר הוא קרוב מאוד ל-1. ערכו 1.0003.
כלומר, הוא גדול מ-1 ב $0.03\%$.

הסטייה היחסית של התוצאה מן הערך הידוע היא:
\begin{lefteqs}
    \frac{\left|c - c_0\right|}{c_0} = \frac{\left|3.049 - 2.998\right|}{2.998} = 0.017 \approx 1.7\%
\end{lefteqs}

הערך הידוע נמצא היטב בתוך טווח השגיאה $[2.86, 3.24] \cdot 10^8\,m/s$, ולכן המדידה מסכימה עם התיאוריה.
כמו כן השגיאה היחסית ($6.4\%$) קטנה מהסף המקובל של $20-25\%$, ולכן התוצאה משמעותית.

בשלב הכנת המערכת נשאלנו כיצד ניתן למדוד את הפרש המופע
בעזרת אוסילוסקופ. האות המשודר והאות הנקלט מחוברים לשני ערוצי האוסילוסקופ ומוצגים בו-זמנית.
שניהם אותות מחזוריים באותו תדר $f = 50MHz$, כלומר בעלי מחזור $T = \frac{1}{f} = 20ns$,
ובין החזיתות המתאימות שלהם נמדדת הסטה בזמן $\Delta t$. הפרש המופע מתקבל מהיחס בין ההסטה למחזור:
\begin{lefteqs}
    \Delta \varphi = 360^\circ \cdot \frac{\Delta t}{T} = 360^\circ \cdot \Delta t \cdot f
\end{lefteqs}
בפועל, כאשר הזזנו את המראה ראינו על המסך את ההסטה בין שני האותות גדלה.
את מדידת הפרש המופע יכלנו לעשות בעזרת ספירת השנתות בין תחילת המופע של הגל היוצא והגל המוחזר, והצבה בנוסחא לעיל.

\newpage
\section*{ניסוי 2 - מדידת מהירות האור באוויר (שיטה 2)}

\subsection*{מטרת הניסוי}
מדידת מהירות האור באוויר מתוך שיפוע הגרף של הפרש הדרכים כפונקציה של הפרש המופע.

\subsection*{רקע תיאורטי}
כפי שהוסבר בניסוי 1, הפרש המופע שצוברת קרן לאורך הפרש דרכים $\Delta x$ הוא:
\begin{equation}
    \Delta \varphi = 2\pi \frac{\Delta x}{\lambda}
\end{equation}
כאשר $\Delta \varphi$ המופע ברדיאנים, $\Delta x$ הפרש הדרכים ביחידות $m$ ו-$\lambda$ אורך הגל של האות המאפנן ביחידות $m$.

עם $\lambda = c/f$, ובמעבר לתצוגת מופע במעלות, מתקבלת נוסחת המדידה:
\begin{equation}
    c = 360 \cdot \frac{\Delta x}{\Delta \varphi} \cdot f
\end{equation}
כאשר $c$ מהירות האור ביחידות $m/s$, $f$ תדר האות המאפנן ביחידות $Hz$ ו-$\Delta \varphi$ הפרש המופע במעלות.

בניסוי 1 נמדדה נקודה בודדת $\Delta \varphi = 180^\circ$. כאן נשנה את $\Delta x$ באופן רציף ונמדוד את
$\Delta \varphi$ המתאים. לפי (2) מתקיים בין השניים קשר ליניארי:
\begin{equation}
    \Delta x = a \cdot \Delta \varphi \,, \qquad a = \frac{c}{360 f}
\end{equation}
כאשר $a$ שיפוע הגרף ביחידות $m/deg$. מבידוד $c$:
\begin{equation}
    c = 360 \cdot a \cdot f
\end{equation}

היתרון של שיטה זו הוא שהשיפוע אינו מושפע מהיסט קבוע בתצוגת המופע, ושהוא מנצל את כל נקודות המדידה
במקום נקודה אחת.

גם כאן הקרן עוברת את המרחק הלוך ושוב, ולכן הפרש הדרכים הוא כפול ממרחק ההזזה של המראה:
\begin{equation}
    \Delta x = 2d
\end{equation}
כאשר $d$ מרחק ההזזה של המראה על המסילה.

\subsection*{מערכת המדידה ומהלך הניסוי}

המערכת נשארה כפי שהייתה בניסוי 1: מד מהירות האור והמראה מורכבים על המסילה, והמכשיר במצב תצוגת מופע במעלות.

מיקמנו את המראה קרוב למד המהירות ואפסנו את תצוגת המופע באמצעות כפתור \begin{english}"Calibration"\end{english}.

לאחר מכן הרחקנו את המראה לאורך המסילה בקפיצות של $10cm$, עשר פעמים, ובכל מיקום רשמנו את קריאת
הפרש המופע שהציג המכשיר.

\subsection*{תוצאות וניתוח}

בטבלה שלהלן $\Delta x$ הוא כבר הפרש הדרכים של הקרן, כלומר כפול ממרחק ההזזה של המראה, לפי (5):

\begin{table}[H]
    \centering
    \renewcommand{\arraystretch}{1.5}
    \begin{tabular}{|c|c|}
        \hline
        \text{[cm] $\Delta x$} & \text{[מעלות] $\Delta \varphi$} \\ \hline
        0 & 0 \\ \hline
        20 & 11 \\ \hline
        40 & 24 \\ \hline
        60 & 36 \\ \hline
        80 & 48 \\ \hline
        100 & 60 \\ \hline
        120 & 73 \\ \hline
        140 & 85 \\ \hline
        160 & 98 \\ \hline
        180 & 109 \\ \hline
        200 & 121 \\ \hline
    \end{tabular}
    \caption{הפרש הדרכים כפונקציה של הפרש המופע}
    \label{tab:l6_e2}
\end{table}

% \begin{figure}[H]
%     \centering
%     \includegraphics[width=0.8\textwidth]{E2_graph.jpg}
%     \caption{גרף של $\Delta x$ כפונקציה של $\Delta \varphi$ עם קו מגמה ליניארי}
% \end{figure}

מהתאמת קו מגמה ליניארי $\Delta x = a \cdot \Delta \varphi + b$ למדידות התקבל:
\begin{lefteqs}
    a = 1.639\,\frac{cm}{deg} \,, \quad b = 0.91\,cm \,, \quad R^2 = 0.9999
\end{lefteqs}

את שגיאות ההתאמה חישבנו לפי נוסחה (17) בתדריך השגיאות:
\begin{lefteqs}
    \Delta y = \sqrt{\frac{1}{N-2} \sum \left(y_i - a x_i - b\right)^2} \\
    \Delta a = \Delta y \sqrt{\frac{N}{A}} \\
    \Delta b = \Delta y \sqrt{\frac{S_{xx}}{A}} \\
    A = N S_{xx} - S_x^2 \,, \quad S_{xx} = \sum x_i^2 \,, \quad S_x = \sum x_i
\end{lefteqs}
כאשר $\Delta y$ מאפיין את סטיית הנקודות מקו המגמה, $N = 11$ מספר המדידות ו-$A$ מבטא את פיזור ערכי $x$.
הערכים חושבו באמצעות פונקציית \begin{english}LINEST\end{english} באקסל המיישמת נוסחה זו:
\begin{lefteqs}
    \Delta y = 0.83\,cm \\
    \Delta a = 0.0065\,\frac{cm}{deg} \\
    \Delta b = 0.47\,cm
\end{lefteqs}

כלומר, בעיגול בהתאם לגודל השגיאה:
\begin{lefteqs}
    a = 1.639 \pm 0.007\,\frac{cm}{deg} \\
    b = 0.9 \pm 0.5\,cm
\end{lefteqs}

יש לציין כי נוסחה (17) מניחה $\Delta x \ll \Delta y$, ובמקרה שלנו דווקא השגיאה בציר $x$ (קריאת המופע,
$1^\circ$, השקולה לכ-$1.6cm$ בציר $y$) גדולה מהשגיאה בקריאת הסרגל. כדי לוודא שהדבר אינו משנה את
התוצאה חישבנו גם את הרגרסיה ההפוכה, של $\Delta \varphi$ כפונקציה של $\Delta x$, וקיבלנו שיפוע
ששקול ל-$1.6393\,cm/deg$ - הפרש של פחות מ-$0.02\%$ מהשיפוע שהתקבל בהתאמה הישירה.
מכאן שהמסקנות אינן תלויות בבחירת הצירים.

כעת נחשב את מהירות האור מהשיפוע לפי (4):
\begin{lefteqs}
    c = 360 \cdot a \cdot f = 360 \cdot 0.01639\,\frac{m}{deg} \cdot 50 \cdot 10^6\,Hz = 2.950 \cdot 10^8\,\frac{m}{s}
\end{lefteqs}

\textbf{שגיאת מהירות האור.} $c$ הוא מכפלה של השיפוע והתדר, ולכן השגיאה היחסית היא סכום השגיאות
היחסיות (נוסחה (8) בתדריך השגיאות):
\begin{lefteqs}
    \frac{\delta c}{c} = \frac{\Delta a}{a} + \frac{\delta f}{f}
\end{lefteqs}
נציב:
\begin{lefteqs}
    \frac{\delta c}{c} = \frac{0.0065}{1.639} + 0.05 = 0.0040 + 0.05 = 0.054
\end{lefteqs}
כלומר שגיאה יחסית של $5.4\%$, ולכן:
\begin{lefteqs}
    \delta c = 0.054 \cdot 2.950 \cdot 10^8 = 0.16 \cdot 10^8\,\frac{m}{s}
\end{lefteqs}

התוצאה הסופית:
\begin{lefteqs}
    c = (2.95 \pm 0.16) \cdot 10^8\,\frac{m}{s}
\end{lefteqs}

\subsection*{דיון ומסקנות}

התוצאות שהתקבלו:
\begin{lefteqs}
    a = 1.639 \pm 0.007\,\frac{cm}{deg} \\
    b = 0.9 \pm 0.5\,cm \\
    c = (2.95 \pm 0.16) \cdot 10^8\,\frac{m}{s}
\end{lefteqs}

\textbf{ליניאריות.} הגרף ליניארי במידה טובה מאוד: מקדם המתאם הוא $R^2 = 0.9999$ והשאריות
(שגודלן האופייני $0.83cm$ מתוך תחום של $200cm$) מפוזרות ללא מגמה שיטתית.
זוהי אישוש ישיר של נוסחה (2), כלומר של הטענה שהמופע שצובר הגל פרופורציוני למרחק שעבר.

החיתוך שהתקבל, $b = 0.9 \pm 0.5\,cm$, אמור להיות אפס: לפני תחילת המדידה איפסנו את התצוגה,
ולכן $\Delta \varphi = 0$ צריך להתאים ל-$\Delta x = 0$. הערך 0 נמצא מחוץ לטווח השגיאה אך במרחק
של פחות משתי שגיאות ממנו, ולכן אינו מעיד על בעיה ממשית. אם קיים כאן היסט אמיתי, הוא שקול
להיסט של כ-$0.5^\circ$ בתצוגת המופע - כלומר בגודל של רזולוציית התצוגה עצמה.

\textbf{השוואה לערך הידוע.} הערך הידוע הוא $c_0 = 2.998 \cdot 10^8\,m/s$. הסטייה היחסית היא:
\begin{lefteqs}
    \frac{\left|c - c_0\right|}{c_0} = \frac{\left|2.950 - 2.998\right|}{2.998} = 0.016 \approx 1.6\%
\end{lefteqs}
הערך הידוע נמצא בתוך טווח השגיאה $[2.79, 3.11] \cdot 10^8\,m/s$, ולכן המדידה מסכימה עם התיאוריה.

\textbf{השוואה לניסוי 1.} בניסוי 1 התקבל $c = (3.05 \pm 0.19) \cdot 10^8\,m/s$ ובניסוי 2
$c = (2.95 \pm 0.16) \cdot 10^8\,m/s$. טווחי השגיאה של שתי התוצאות חופפים ברובם, ושתיהן מסכימות עם
הערך הידוע.

עם זאת יש לשים לב לנקודה מהותית בהשוואה: אי-הוודאות בתדר ($5\%$) משותפת לשתי המדידות, שכן שתיהן
השתמשו באותו מכשיר ובאותו נתון יצרן. לכן היא אינה מבדילה ביניהן ומבטלת עצמה בהשוואה.
המרכיבים הבלתי תלויים בלבד הם $1.4\%$ בניסוי 1 ו-$0.4\%$ בניסוי 2, בעוד שההפרש בפועל בין שתי
התוצאות הוא $3.3\%$ - כלומר גדול פי כמה מהשגיאות האקראיות שאינן משותפות.
מכאן שקיים הפרש שיטתי קטן בין שתי השיטות שאיננו מוסבר על ידי השגיאות שהערכנו.

ההסבר הסביר ביותר הוא ההבדל באופן שבו נקראת התצוגה. ניסוי 1 מסתמך על נקודת מופע יחידה
($180^\circ$) המתקבלת במרחק של כמטר וחצי, שם עוצמת האות החוזר נמוכה והקריאה פחות יציבה,
וההכרעה מתי "הופיע" הערך $180^\circ$ נעשית בתוך תחום המתה של $7mm$. ניסוי 2, לעומת זאת,
משקלל אחת-עשרה נקודות לאורך תחום של מטר ומחלץ שיפוע, כך שהיסטים קבועים בקריאה מתבטלים
והשפעת תחום ההמתה ממוצעת. מסיבות אלה יש להעדיף את תוצאת ניסוי 2 כמדויקת יותר,
וזה אכן משתקף בשגיאה האקראית הקטנה יותר שלה.

\textbf{מקורות שגיאה ושיפורים.} כמו בניסוי 1, השגיאה נשלטת כמעט לחלוטין על ידי אי-הוודאות בתדר
($5\%$ מתוך $5.4\%$), ותרומת ההתאמה עצמה זניחה בהשוואה. השיפור המשמעותי היחיד הוא מדידה ישירה
של תדר האות המאפנן. שיפור משני הוא הארכת המסילה: השיפוע נקבע טוב יותר ככל שתחום $\Delta x$ גדול יותר.

\newpage
\section*{ניסוי 3 - מדידת מקדם השבירה ומהירות האור במים}

\subsection*{מטרת הניסוי}
מדידת מקדם השבירה של מים וחישוב מהירות האור בהם, מתוך מרחק ההזזה של המראה הדרוש כדי לפצות על
העיכוב שגורם עמוד המים במסלול הקרן.

\subsection*{רקע תיאורטי}
בסביבה חומרית שקופה מהירות הגל האלקטרומגנטי קטנה ממהירותו בריק. היחס בין השתיים הוא מקדם השבירה:
\begin{equation}
    n = \frac{c}{c_m}
\end{equation}
כאשר $c$ מהירות האור בריק (ובקירוב טוב גם באוויר) ביחידות $m/s$, ו-$c_m$ מהירותו בחומר ביחידות $m/s$.
מהגדרה זו נובע תמיד $n \geq 1$.

התדירות של הגל אינה משתנה במעבר בין תווכים, ולכן קיצור המהירות מתבטא בקיצור אורך הגל:
\begin{equation}
    \lambda_m = \frac{c_m}{f} = \frac{\lambda}{n}
\end{equation}
כאשר $\lambda$ אורך הגל בריק ו-$\lambda_m$ אורך הגל בחומר, שניהם ביחידות $m$.

כפי שראינו, הפרש המופע שצוברת קרן לאורך מרחק $\Delta x$ בתווך שאורך הגל בו הוא $\lambda$ הוא
$\Delta \varphi = 2\pi \Delta x / \lambda$. מכיוון ש-$\lambda_m$ קטן מ-$\lambda$, אותו מרחק גיאומטרי
בתוך החומר "עולה" יותר מופע - וזה מה שהניסוי מודד.

נסמן ב-$l_m$ את אורך גליל החומר וב-$l_1$ את המרחק מהמכשיר למראה במדידה הראשונה.
במדידה הראשונה הגליל נמצא במסלול הקרן, ולכן הקרן עוברת הלוך ושוב מרחק $2l_m$ בתוך החומר
ומרחק $2(l_1 - l_m)$ באוויר. המופע שהיא צוברת הוא:
\begin{equation}
    \Delta \varphi_1 = 2\pi \frac{2 l_m}{\lambda_m} + 2\pi \frac{2(l_1 - l_m)}{\lambda}
    = \frac{4\pi}{\lambda}\left( n l_m - l_m + l_1 \right)
\end{equation}

במדידה השנייה מסירים את הגליל ומזיזים את המראה מרחק $d$ עד שתצוגת המופע חוזרת לאפס.
כעת כל המסלול הוא באוויר:
\begin{equation}
    \Delta \varphi_2 = 2\pi \frac{2 l_2}{\lambda} = \frac{4\pi}{\lambda}\left( l_1 + d \right)
\end{equation}
כאשר $l_2 = l_1 + d$ המרחק החדש למראה.

מכיוון שבשתי המדידות תצוגת המופע מראה אפס, המופע שצברה הקרן זהה, כלומר $\Delta \varphi_1 = \Delta \varphi_2$:
\begin{lefteqs}
    n l_m - l_m + l_1 = l_1 + d
\end{lefteqs}
ומכאן מקדם השבירה:
\begin{equation}
    n = 1 + \frac{d}{l_m}
\end{equation}
כאשר $d$ מרחק הזזת המראה ביחידות $cm$ ו-$l_m$ אורך גליל החומר ביחידות $cm$.

הפירוש הפיזיקלי פשוט: $d$ הוא אורך הדרך הנוספת באוויר ששקולה לעיכוב שגורם החומר,
והיחס $d / l_m$ הוא בדיוק העיכוב היחסי, $n - 1$.

לבסוף, מהירות האור בחומר מתקבלת מ-(1):
\begin{equation}
    c_m = \frac{c}{n}
\end{equation}

\subsection*{מערכת המדידה ומהלך הניסוי}

מדדנו בסרגל את אורך הגליל המלא במים וקיבלנו $l_m = 52.5cm$, בדיוק של $1mm$.

הרכבנו את הגליל על המסילה בשני מחזיקים, בין מד מהירות האור לבין המראה, כך שקרן האור עוברת דרכו לאורך צירו.

אפסנו את תצוגת המופע באמצעות כפתור \begin{english}"Calibration"\end{english} ורשמנו את קריאת מיקום המראה.

לאחר מכן הסרנו את הגליל מהמסילה והזזנו את המראה עד שתצוגת המופע חזרה להראות אפס, ורשמנו שוב את קריאת המיקום.

חזרנו על התהליך חמש פעמים בסך הכול, כאשר בכל פעם מיקום המראה ההתחלתי היה שונה.

\subsection*{תוצאות וניתוח}

תוצאות המדידה, כאשר $d$ הוא הפרש קריאות המיקום:

\begin{table}[H]
    \centering
    \renewcommand{\arraystretch}{1.5}
    \begin{tabular}{|c|c|}
        \hline
        \text{[cm] $d$} & \text{[cm] מיקום התחלתי} \\ \hline
        15.75 & 60 \\ \hline
        15.85 & 70 \\ \hline
        15.85 & 75 \\ \hline
        13.85 & 80 \\ \hline
        17.00 & 65 \\ \hline
    \end{tabular}
    \caption{מרחק הזזת המראה לאיפוס המופע לאחר הסרת גליל המים}
    \label{tab:l6_e3}
\end{table}

הערך הטוב ביותר עבור $d$ הוא ממוצע חמש המדידות (נוסחה (1) בתדריך השגיאות):
\begin{lefteqs}
    \bar{d} = \frac{1}{N}\sum_{i=1}^{N} d_i = \frac{15.75 + 15.85 + 15.85 + 13.85 + 17.00}{5} = 15.66\,cm
\end{lefteqs}

\textbf{שגיאה סטטיסטית} (נוסחה (2) בתדריך השגיאות):
\begin{lefteqs}
    \Delta_s d = \sqrt{\frac{\sum_{i=1}^{N}\left(d_i - \bar{d}\right)^2}{N}} \\
    \Delta_s d = \sqrt{\frac{0.09^2 + 0.19^2 + 0.19^2 + (-1.81)^2 + 1.34^2}{5}} = 1.02\,cm
\end{lefteqs}

\textbf{שגיאת מכשיר.} כמו בניסוי 1, השגיאה הדומיננטית היא תחום ההמתה של תצוגת המופע.
המדידה כאן היא מדידת אפס: מזיזים את המראה עד שהתצוגה חוזרת ל-$0$, ולכן אותו תחום המתה של $7mm$
חל על שתי קריאות המיקום, ההתחלתית והסופית:
\begin{lefteqs}
    \Delta_m d = 0.35 + 0.35 = 0.7\,cm
\end{lefteqs}

\textbf{שגיאה כוללת} (נוסחה (4) בתדריך השגיאות):
\begin{lefteqs}
    \Delta d = \sqrt{\left(\Delta_m d\right)^2 + \left(\Delta_s d\right)^2} = \sqrt{0.7^2 + 1.02^2} = 1.23\,cm
\end{lefteqs}

כעת נחשב את מקדם השבירה לפי (5):
\begin{lefteqs}
    n = 1 + \frac{\bar{d}}{l_m} = 1 + \frac{15.66}{52.5} = 1.298
\end{lefteqs}

\textbf{שגיאת מקדם השבירה.} $n$ הוא פונקציה של שני גדלים נמדדים, $d$ ו-$l_m$.
נשתמש בשיטת הדיפרנציאל (נוסחה (11) בתדריך השגיאות):
\begin{lefteqs}
    \Delta n = \left| \frac{\partial n}{\partial d} \right| \Delta d + \left| \frac{\partial n}{\partial l_m} \right| \Delta l_m
    = \frac{\Delta d}{l_m} + \frac{\bar{d}}{l_m^2} \Delta l_m
\end{lefteqs}
נציב:
\begin{lefteqs}
    \Delta n = \frac{1.23}{52.5} + \frac{15.66}{52.5^2} \cdot 0.1 = 0.0235 + 0.0006 = 0.024
\end{lefteqs}
כלומר תרומת מדידת אורך הגליל זניחה לחלוטין ביחס לתרומת $d$. התוצאה:
\begin{lefteqs}
    n = 1.30 \pm 0.02
\end{lefteqs}

\textbf{מהירות האור במים.} לפי (6), עם הערך הידוע $c_0 = 2.998 \cdot 10^8\,m/s$:
\begin{lefteqs}
    c_m = \frac{c_0}{n} = \frac{2.998 \cdot 10^8}{1.298} = 2.309 \cdot 10^8\,\frac{m}{s}
\end{lefteqs}
בחרנו להשתמש בערך הידוע ולא בתוצאות ניסויים 1-2, מכיוון שאלה נושאות שגיאה של כ-$5\%$ הנובעת
כמעט כולה מאי-הוודאות בתדר, ושימוש בהן היה מטשטש את הדיוק של מדידת $n$ עצמה, שאינה תלויה בתדר כלל.

השגיאה מתקבלת שוב מהדיפרנציאל:
\begin{lefteqs}
    \Delta c_m = \left| \frac{\partial c_m}{\partial n} \right| \Delta n = \frac{c_0}{n^2} \Delta n
    = \frac{2.998 \cdot 10^8}{1.298^2} \cdot 0.024 = 0.04 \cdot 10^8\,\frac{m}{s}
\end{lefteqs}

התוצאה הסופית:
\begin{lefteqs}
    n = 1.30 \pm 0.02 \\
    c_m = (2.31 \pm 0.04) \cdot 10^8\,\frac{m}{s}
\end{lefteqs}

\subsection*{דיון ומסקנות}

התוצאות שהתקבלו עבור מים:
\begin{lefteqs}
    n = 1.30 \pm 0.02 \\
    c_m = (2.31 \pm 0.04) \cdot 10^8\,\frac{m}{s}
\end{lefteqs}

הערך הידוע מן הספרות למקדם השבירה של מים באור נראה הוא $n_{lit} = 1.333$
(מים בטמפרטורת החדר, אורך גל $589nm$). מכאן מהירות האור הצפויה במים היא
$c_0 / n_{lit} = 2.249 \cdot 10^8\,m/s$. ערך זה עקבי עם הנתון המופיע ברקע ההיסטורי של הניסוי,
לפיו מהירות האור במים נמדדה כקטנה בערך בשלושה רבעים ממהירותה בריק.

הסטיות היחסיות הן:
\begin{lefteqs}
    \frac{\left| n - n_{lit} \right|}{n_{lit}} = \frac{\left| 1.298 - 1.333 \right|}{1.333} = 0.026 \approx 2.6\% \\
    \frac{\left| c_m - c_0/n_{lit} \right|}{c_0/n_{lit}} = \frac{\left| 2.309 - 2.249 \right|}{2.249} = 0.027 \approx 2.7\%
\end{lefteqs}

הערך $1.333$ נמצא ממש בקצה טווח השגיאה $[1.274, 1.322]$ ומעט מחוצה לו, כלומר המדידה מסכימה עם
הערך הידוע רק בגבול. השגיאה היחסית ב-$n-1$, כלומר בגודל שנמדד בפועל, היא $\Delta d / \bar{d} = 7.9\%$,
ולכן התוצאה משמעותית, אך היא הפחות מדויקת מבין ארבעת הניסויים.

\textbf{הערה על סוג מקדם השבירה.} יש לציין כי השיטה כאן אינה מודדת את מקדם השבירה במובן הרגיל
(היחס בין מהירויות המופע), אלא את היחס בין מהירויות המעטפת: מה שמתעכב בחומר הוא אות האפנון
בתדר $50MHz$ הרכוב על גל האור, ולא הגל האופטי עצמו. מבחינה טכנית זהו מקדם השבירה של החבורה
(\begin{english}group index\end{english}), שבמים גדול במעט מזה של המופע.
משמעות הדבר היא שההפרש בין התוצאה שלנו לערך הספרותי אינו נובע מהבחנה זו - אילו התחשבנו בה,
הערך הצפוי היה גדל והסטייה הייתה גדלה, לא קטנה.

\textbf{מקורות השגיאה.} כפי שראינו, השגיאה נובעת כמעט כולה ממדידת $d$, ובתוכה המרכיב הסטטיסטי
($1.02cm$) גדול ממרכיב המכשיר ($0.7cm$). הפיזור בין המדידות גדול במיוחד כאן:
המדידה השלישית והחמישית ($13.85cm$ ו-$17.00cm$) רחוקות מהממוצע ביותר משגיאה אחת, ולבדן
אחראיות לרוב הפיזור. לא פסלנו אף מדידה, מכיוון שאין שום קריטריון ניסויי המצדיק זאת.

שלושה גורמים בלתי נפרדים תורמים לפיזור זה, וכולם מחריפים בהשוואה לניסויים 1-2:
ראשית, המרחק הנמדד קצר ($\sim 16cm$ לעומת $\sim 152cm$), ולכן אותו תחום המתה של $7mm$ תופס
נתח גדול פי עשרה מהתוצאה. שנית, הכנסת הגליל והוצאתו משנות את יישור הקרן: פני החומר אינם מאונכים
בדיוק לקרן, וגליל מים פועל גם כעדשה גלילית המשנה את התכנסות האלומה ואת עוצמת האות המגיע לגלאי.
שלישית, האורך $l_m = 52.5cm$ הוא האורך החיצוני של הגליל, הכולל את חלונות הקצה, ואילו עמוד המים עצמו
קצר במקצת ממנו - כך שהגודל שחילקנו בו אינו בדיוק הגודל שמופיע בגזירה של (5).

\textbf{שיפורים.} השיפור הישיר ביותר הוא הגדלת $d$ ביחס לתחום ההמתה, על ידי שימוש בגליל ארוך יותר
או במספר גלילים בטור. שיפור נוסף הוא מדידה נפרדת של אורך עמוד המים ושל עובי חלונות הקצה,
כדי להפריד את תרומתם. לבסוף, הגדלת מספר החזרות תקטין את המרכיב הסטטיסטי, וביצוע המדידה גם בכיוון
ההפוך (הכנסת הגליל במקום הוצאתו) יאפשר לזהות הטיה שיטתית הנובעת מכיוון הזזת המראה.

\newpage
\section*{ניסוי 4 - מדידת מקדם השבירה ומהירות האור בפלסטיק אקרילי}

\subsection*{מטרת הניסוי}
מדידת מקדם השבירה של פלסטיק אקרילי וחישוב מהירות האור בו, באותה שיטה שבה נמדדו המים בניסוי 3.

\subsection*{רקע תיאורטי}
הרקע התיאורטי זהה לזה של ניסוי 3, ונביא כאן בקצרה את הנוסחאות הדרושות.

מקדם השבירה של חומר מוגדר כיחס בין מהירות האור בריק (ובקירוב באוויר) לבין מהירותו בחומר:
\begin{equation}
    n = \frac{c}{c_m}
\end{equation}
כאשר $c$ מהירות האור בריק ו-$c_m$ מהירות האור בחומר, שתיהן ביחידות $m/s$.

מכיוון שאורך הגל בחומר קצר בפקטור $n$, קטע חומר באורך $l_m$ גורם לקרן לצבור מופע נוסף
לעומת אותו קטע באוויר. אם לאחר הסרת החומר יש להזיז את המראה מרחק $d$ כדי להחזיר את
הפרש המופע לאפס, אזי הדרך הנוספת באוויר שקולה לעיכוב שגרם החומר, ומתקבל:
\begin{equation}
    n = 1 + \frac{d}{l_m}
\end{equation}
כאשר $d$ מרחק הזזת המראה ו-$l_m$ אורך גליל החומר, שניהם ביחידות $cm$.

מהירות האור בחומר מתקבלת מ-(1):
\begin{equation}
    c_m = \frac{c}{n}
\end{equation}

\subsection*{מערכת המדידה ומהלך הניסוי}

החלפנו את מכל המים בגליל של פלסטיק אקרילי ומדדנו בסרגל את אורכו: $l_m = 49cm$.

הרכבנו את הגליל על המסילה בין מד מהירות האור למראה, וחזרנו על אותם שלבים שביצענו בניסוי 3:
איפסנו את תצוגת המופע, רשמנו את קריאת מיקום המראה, הסרנו את הגליל, הזזנו את המראה עד שהמופע חזר
לאפס ורשמנו שוב את קריאת המיקום.

גם כאן חזרנו על התהליך חמש פעמים בסך הכול, כאשר בכל פעם מיקום המראה ההתחלתי היה שונה.

\subsection*{תוצאות וניתוח}

תוצאות המדידה, כאשר $d$ הוא הפרש קריאות המיקום:

\begin{table}[H]
    \centering
    \renewcommand{\arraystretch}{1.5}
    \begin{tabular}{|c|c|}
        \hline
        \text{[cm] $d$} & \text{[cm] מיקום התחלתי} \\ \hline
        23.0 & 50 \\ \hline
        22.3 & 55 \\ \hline
        23.0 & 60 \\ \hline
        22.3 & 65 \\ \hline
        21.5 & 70 \\ \hline
    \end{tabular}
    \caption{מרחק הזזת המראה לאיפוס המופע לאחר הסרת גליל האקריל}
    \label{tab:l6_e4}
\end{table}

הערך הטוב ביותר עבור $d$ הוא ממוצע חמש המדידות (נוסחה (1) בתדריך השגיאות):
\begin{lefteqs}
    \bar{d} = \frac{1}{N}\sum_{i=1}^{N} d_i = \frac{23.0 + 22.3 + 23.0 + 22.3 + 21.5}{5} = 22.42\,cm
\end{lefteqs}

\textbf{שגיאה סטטיסטית} (נוסחה (2) בתדריך השגיאות):
\begin{lefteqs}
    \Delta_s d = \sqrt{\frac{\sum_{i=1}^{N}\left(d_i - \bar{d}\right)^2}{N}} \\
    \Delta_s d = \sqrt{\frac{0.58^2 + (-0.12)^2 + 0.58^2 + (-0.12)^2 + (-0.92)^2}{5}} = 0.56\,cm
\end{lefteqs}

\textbf{שגיאת מכשיר.} כמו בניסוי 3, מדובר במדידת אפס שבה תחום ההמתה של $7mm$ חל על שתי קריאות המיקום:
\begin{lefteqs}
    \Delta_m d = 0.35 + 0.35 = 0.7\,cm
\end{lefteqs}
כאן, בשונה מניסוי 3, מרכיב המכשיר גדול מהמרכיב הסטטיסטי.

\textbf{שגיאה כוללת} (נוסחה (4) בתדריך השגיאות):
\begin{lefteqs}
    \Delta d = \sqrt{\left(\Delta_m d\right)^2 + \left(\Delta_s d\right)^2} = \sqrt{0.7^2 + 0.56^2} = 0.89\,cm
\end{lefteqs}

מקדם השבירה לפי (2):
\begin{lefteqs}
    n = 1 + \frac{\bar{d}}{l_m} = 1 + \frac{22.42}{49} = 1.458
\end{lefteqs}

\textbf{שגיאת מקדם השבירה} (שיטת הדיפרנציאל, נוסחה (11) בתדריך השגיאות):
\begin{lefteqs}
    \Delta n = \left| \frac{\partial n}{\partial d} \right| \Delta d + \left| \frac{\partial n}{\partial l_m} \right| \Delta l_m
    = \frac{\Delta d}{l_m} + \frac{\bar{d}}{l_m^2} \Delta l_m
\end{lefteqs}
נציב, כאשר $\Delta l_m = 0.1cm$ (אותו סרגל שבו נמדד גליל המים):
\begin{lefteqs}
    \Delta n = \frac{0.89}{49} + \frac{22.42}{49^2} \cdot 0.1 = 0.0182 + 0.0009 = 0.019
\end{lefteqs}
כלומר:
\begin{lefteqs}
    n = 1.46 \pm 0.02
\end{lefteqs}

\textbf{מהירות האור באקריל} לפי (3), עם הערך הידוע $c_0 = 2.998 \cdot 10^8\,m/s$:
\begin{lefteqs}
    c_m = \frac{c_0}{n} = \frac{2.998 \cdot 10^8}{1.458} = 2.057 \cdot 10^8\,\frac{m}{s} \\
    \Delta c_m = \frac{c_0}{n^2} \Delta n = \frac{2.998 \cdot 10^8}{1.458^2} \cdot 0.019 = 0.03 \cdot 10^8\,\frac{m}{s}
\end{lefteqs}

התוצאה הסופית:
\begin{lefteqs}
    n = 1.46 \pm 0.02 \\
    c_m = (2.06 \pm 0.03) \cdot 10^8\,\frac{m}{s}
\end{lefteqs}

\subsection*{דיון ומסקנות}

התוצאות שהתקבלו עבור פלסטיק אקרילי:
\begin{lefteqs}
    n = 1.46 \pm 0.02 \\
    c_m = (2.06 \pm 0.03) \cdot 10^8\,\frac{m}{s}
\end{lefteqs}

הערך הידוע מן הספרות למקדם השבירה של פלסטיק אקרילי (\begin{english}PMMA\end{english}) באור נראה הוא
$n_{lit} = 1.49$, ומכאן מהירות האור הצפויה בו היא $c_0 / n_{lit} = 2.012 \cdot 10^8\,m/s$.

הסטיות היחסיות הן:
\begin{lefteqs}
    \frac{\left| n - n_{lit} \right|}{n_{lit}} = \frac{\left| 1.458 - 1.49 \right|}{1.49} = 0.022 \approx 2.2\% \\
    \frac{\left| c_m - c_0/n_{lit} \right|}{c_0/n_{lit}} = \frac{\left| 2.057 - 2.012 \right|}{2.012} = 0.022 \approx 2.2\%
\end{lefteqs}

הערך $1.49$ נמצא בתוך טווח השגיאה $[1.439, 1.477]$, ולכן המדידה מסכימה עם הערך הידוע.
השגיאה היחסית בגודל הנמדד בפועל היא $\Delta d / \bar{d} = 4.0\%$, ולכן התוצאה משמעותית.

\textbf{השוואה לניסוי 3.} מדידת האקריל יצאה מדויקת יותר ממדידת המים, ומשתי סיבות מבניות:
ראשית, מקדם השבירה גבוה יותר ולכן $d$ גדול יותר ($22.4cm$ לעומת $15.7cm$), כך שאותו תחום המתה של
$7mm$ תופס נתח קטן יותר מהתוצאה. שנית, האקריל הוא גוף מוצק אחיד: אין בו חלונות קצה נפרדים ואין בו
עדשה גלילית של נוזל, ולכן הכנסתו והוצאתו משנות את יישור הקרן פחות. הדבר משתקף ישירות בפיזור
המדידות - $0.56cm$ באקריל לעומת $1.02cm$ במים.

\textbf{מקורות השגיאה.} כאן, בשונה מניסוי 3, השגיאה נשלטת על ידי רזולוציית תצוגת המופע ($0.7cm$)
ולא על ידי הפיזור הסטטיסטי ($0.56cm$). המשמעות היא שהגדלת מספר החזרות לא תשפר את התוצאה משמעותית:
המדידה הגיעה לגבול הרזולוציה של המכשיר.

\textbf{שיפורים.} השיפור היעיל הוא הגדלת $d$ ביחס לתחום ההמתה, על ידי שימוש בגליל אקריל ארוך יותר או
בכמה גלילים בטור. בנוסף, ניתן לצמצם את תחום ההמתה עצמו על ידי הגברת עוצמת האות החוזר -
למשל על ידי יישור מדוקדק יותר של המראה ועבודה במרחקים קצרים ככל האפשר.

\textbf{סיכום כללי של המעבדה.} בארבעת הניסויים נמדדו:
\begin{lefteqs}
    \text{Exp. 1:} \quad c = (3.05 \pm 0.19) \cdot 10^8\,\frac{m}{s} \\
    \text{Exp. 2:} \quad c = (2.95 \pm 0.16) \cdot 10^8\,\frac{m}{s} \\
    \text{Exp. 3:} \quad n_{water} = 1.30 \pm 0.02 \,, \quad c_m = (2.31 \pm 0.04) \cdot 10^8\,\frac{m}{s} \\
    \text{Exp. 4:} \quad n_{acrylic} = 1.46 \pm 0.02 \,, \quad c_m = (2.06 \pm 0.03) \cdot 10^8\,\frac{m}{s}
\end{lefteqs}

כל התוצאות מסכימות עם הערכים הידועים מן הספרות בסטייה של $1.6\% - 2.7\%$.
ראוי לשים לב להבדל מבני בין שני חצאי המעבדה: מדידת $c$ עצמה מוגבלת כמעט לחלוטין על ידי אי-הוודאות
בתדר האות המאפנן ($5\%$), נתון יצרן שאין לנו שליטה עליו, בעוד שמדידת $n$ אינה תלויה בתדר כלל
ומוגבלת על ידי רזולוציית תצוגת המופע. לכן, למרות שמדידות $n$ נראות "עקיפות" יותר, הן דווקא
המדויקות מבין השתיים.

\end{document}